% Pursuing Stacks (cotes 134-2 to 134-8): archive cote 134-6, site page 81 (PDF page 82), typescript page 399.
% typescript. Draft compared: ps5.tex:5584-5671.
% AI-transcribed by Claude Opus 5.5 (claude-opus-5-5, adaptive thinking, effort high), Message Batches API, batch ps3 (msgbatch_01XdFmuVVBAcgc7EFxdtCqsj), 2026-11-12; not yet checked by eye.
\scanpage{134-6/81}
\mspage{399}

(49) \qquad $B^{\wedge}_{ab} \hand{\xrightarrow{\approx}} \underline{Hom}_{\hand{!}}(A^{\wedge}_{ab} , (Ab))$ ,

which should be viewed as the "abelian" analogon of the equivalence (41) \xout{\unsure{al}}
above (corresponding to the case \hand{$\mathcal{M}$}= (Sets) ). This equivalence again
may be viewed as described (in analogy to (42)) by a canonical pairing

(50) \qquad $A^{\wedge}_{ab} \text{ x } B^{\wedge}_{ab} \longrightarrow (Ab)$ ,

which commutes to small direct limits in each variable i.e. can be viewed
as an object in the category of"abelian bi-cosheaves"

(*) \qquad $\underline{Hom}_{\hand{!!}}(A^{\wedge}_{ab},B^{\wedge}_{ab}; (Ab))$ ,

and gives rise simultaneously to the equivalence (49), and to the symme-%
tric equivalence

(49') \qquad $A^{\wedge}_{ab} \hand{\xrightarrow{\approx}} \underline{Hom}_{\hand{!}}(B^{\wedge}_{ab}, (Ab) )$

of the category of abelian sheaves on $A^{\wedge}$\hand{/}with the category of abelian
cosheaves on $B^{\wedge}$ (which is just (49) with A replaced by B, up to c anoni-%
cal isomorphism at any rate).

Using the equivalence

\hand{(51)} \qquad $\underline{Hom}_{\hand{!!}}(A^{\wedge}_{ab},B^{\wedge}_{ab} ; \underset{\hand{\smile}}{(Ab)}) \hand{\xrightarrow{\approx}} \underline{Hom}(A\text{x}B, \underset{\hand{\smile}}{(Ab)})$

(which is a particular case of the evident abelian analogon of the equi-%
valence (44)), we see that the pairing (50) is described, up to canoni-%
cal isomorphism, by it's \xout{restrictionx} composition with
\[(a,b)\hand{\rightarrow}(Wh_{A^{\wedge}}(a), Wh_{B^{\wedge}}(b)) : A\text{x}B \hand{\longrightarrow} \deleted{(}A^{\wedge}_{ab} \text{ x } B^{\wedge}_{ab} \quad ,\]
and the latter, as is readily ch\hand{\unsure{ec}}ked, is given by

(5\over{1}{2}) \qquad $Wh_A(a) *_{\underline{Z}} Wh_B(b^o) \hand{\simeq} \underline{Z}^{(Hom(b,a))}$ \qquad for a,b in A ,

where we have written $Wh_A$ instead of $Wh_{A^{\wedge}}$ for brevity and accordingly for
B , and where the pairing (50) is denot ed by the symbol $*_{\underline{Z}}$ , in analo-%
gy with the notation $*$ in (43), the index $\underline{Z}$ being added in order to
avoid confusi on with the non-abelian case (43) (and the index being drop-%
ped when no such confusion is to be feared). T\over{h}{\unsure{h}}e formula (5\over{1}{2}) can be
written, with different notations

(5\over{1}{2}')\xout{xx} \qquad $\underline{Z}^{(a)} *_{\underline{Z}} \underline{Z}^{(b^o)} \hand{\simeq} \underline{Z}^{(Hom(b,a))}$ \qquad for a,b in A .

Comparing with the similar formula (45), this suggests the generalization

(5\over{2}{3}) \qquad $Wh_A(F) *_{\underline{Z}} Wh_B(G) \hand{\simeq} \underline{Z}^{(F*G)}$ \qquad ,

or with the exponential notation

(5\over{2}{3}') \qquad $\underline{Z}^{(F)} *_{\underline{Z}} \underline{Z}^{(G)} \hand{\simeq} \underline{Z}^{\hand{(}F*G)}$ \qquad ,

\margin{81}
